%=======================================================================
%  CHE F313 - Separation Processes II
%  PROBLEMS AND SOLUTIONS  (questions from the uploaded slide decks)
%
%  PROVENANCE DISCIPLINE USED THROUGHOUT THIS FILE
%    [SOURCE]  transcribed from the uploaded decks in ../Slides/ (trusted)
%    [TEXT]    from the cited textbook (secondary confirmation only)
%    [DERIVED] computed here - arithmetic/algebra generated, never a source
%    [VERIFY]  independent check of a stated answer
%    [FLAG]    ambiguity, missing datum, or assumption made - read these
%
%  Compile:  tectonic SP2-problems-and-solutions.tex   (or latexmk -xelatex)
%=======================================================================
\documentclass[10pt,a4paper]{article}

\usepackage[a4paper,left=18mm,right=18mm,top=15mm,bottom=18mm]{geometry}
\usepackage[T1]{fontenc}
\usepackage{amsmath,amssymb,mathtools}
\usepackage{xcolor}
\usepackage{booktabs,tabularx,array,longtable}
\usepackage{enumitem}
\usepackage{fancyhdr}
\usepackage{tcolorbox}
\tcbuselibrary{breakable,skins}
\usepackage{siunitx}
\usepackage[colorlinks=true,linkcolor=accent,urlcolor=accent2]{hyperref}

\definecolor{accent}{HTML}{0B3C5D}
\definecolor{accent2}{HTML}{1D6FA5}
\definecolor{srccol}{HTML}{1B5E20}      % SOURCE  - green
\definecolor{srcbg}{HTML}{E8F5E9}
\definecolor{devcol}{HTML}{0D47A1}      % DERIVED - blue
\definecolor{devbg}{HTML}{E3F2FD}
\definecolor{vercol}{HTML}{4A148C}      % VERIFY  - purple
\definecolor{verbg}{HTML}{F3E5F5}
\definecolor{flagcol}{HTML}{B71C1C}     % FLAG    - red
\definecolor{flagbg}{HTML}{FFEBEE}
\definecolor{txtcol}{HTML}{6D4C41}      % TEXT    - brown
\definecolor{txtbg}{HTML}{EFEBE9}
\definecolor{boxbg}{HTML}{F1F4F7}
\definecolor{boxframe}{HTML}{C6D1DA}
\definecolor{rule}{HTML}{9FB3C2}
\definecolor{inbox}{HTML}{0E2430}
\definecolor{notecol}{HTML}{46535C}

\sisetup{detect-all,per-mode=symbol}
\newcommand{\microm}{\ensuremath{\,\mu\mathrm{m}}}

% ---------------------------------------------------------------- chips
\newcommand{\chip}[2]{{\setlength{\fboxsep}{3pt}%
  \colorbox{#1}{\textcolor{white}{\bfseries\footnotesize #2}}}}
\newcommand{\SOURCE}{\chip{srccol}{SOURCE}}
\newcommand{\DERIVED}{\chip{devcol}{DERIVED}}
\newcommand{\VERIFY}{\chip{vercol}{VERIFY}}
\newcommand{\FLAGTAG}{\chip{flagcol}{FLAG}}
\newcommand{\TEXT}{\chip{txtcol}{TEXT}}

\newtcolorbox{srcres}{breakable,enhanced,colback=srcbg,colframe=srccol!60,
  boxrule=0.6pt,arc=1.5pt,left=5pt,right=5pt,top=4pt,bottom=4pt,
  before skip=4pt,after skip=5pt,fontupper=\small}
\newtcolorbox{devres}{breakable,enhanced,colback=devbg,colframe=devcol!55,
  boxrule=0.6pt,arc=1.5pt,left=5pt,right=5pt,top=4pt,bottom=4pt,
  before skip=4pt,after skip=5pt,fontupper=\small}
\newtcolorbox{verres}{breakable,enhanced,colback=verbg,colframe=vercol!55,
  boxrule=0.6pt,arc=1.5pt,left=5pt,right=5pt,top=4pt,bottom=4pt,
  before skip=4pt,after skip=5pt,fontupper=\small}
\newtcolorbox{flagres}{breakable,enhanced,colback=flagbg,colframe=flagcol!60,
  boxrule=0.6pt,arc=1.5pt,left=5pt,right=5pt,top=4pt,bottom=4pt,
  before skip=4pt,after skip=5pt,fontupper=\small}
\newtcolorbox{ansbox}[1][]{enhanced,colback=boxbg,colframe=accent,boxrule=0.9pt,
  arc=1.5pt,left=6pt,right=6pt,top=5pt,bottom=5pt,before skip=5pt,after skip=6pt,
  fontupper=\small\color{inbox},#1}
\newcommand{\notetext}[1]{\par\vspace{1pt}\noindent{\small\color{notecol}#1}\par\vspace{2pt}}

% ---------------------------------------------------------------- headings
\newcommand{\spsection}[1]{%
  \par\addvspace{10pt}%
  \noindent\begin{tcolorbox}[colback=accent,colframe=accent,boxrule=0pt,arc=1.2pt,
      left=7pt,right=7pt,top=5pt,bottom=5pt,width=\linewidth,before skip=0pt,after skip=7pt]
    \color{white}\bfseries\large #1
  \end{tcolorbox}\par\vspace{-2pt}\nopagebreak}
\newcommand{\spsub}[1]{\par\addvspace{7pt}\noindent{\color{accent}\bfseries #1}\par\addvspace{3pt}\nopagebreak}
\newcommand{\qtitle}[2]{%
  \par\addvspace{9pt}\noindent
  \begin{tcolorbox}[colback=accent2!12,colframe=accent2,boxrule=0.8pt,arc=1.5pt,
      left=7pt,right=7pt,top=4pt,bottom=4pt,before skip=0pt,after skip=4pt,width=\linewidth]
    {\large\bfseries\color{accent}#1}\hfill{\small\color{accent2}#2}
  \end{tcolorbox}\par\vspace{1pt}}

\setlist[itemize]{leftmargin=1.1em,topsep=2pt,itemsep=1.5pt,parsep=0pt,label=\textbullet}
\setlist[enumerate]{leftmargin=1.5em,topsep=2pt,itemsep=1.5pt}
\setlength{\parindent}{0pt}
\emergencystretch=2em
\setlength{\parskip}{2.5pt}

\pagestyle{fancy}
\fancyhf{}
\renewcommand{\headrulewidth}{0pt}
\renewcommand{\footrulewidth}{0.4pt}
\fancyfoot[L]{\footnotesize\color{notecol}CHE F313 \textperiodcentered{} Problems \& Solutions \textperiodcentered{} source-tagged}
\fancyfoot[R]{\footnotesize\color{notecol}Page \thepage}

\begin{document}

%=============================================================== cover
\begin{center}
{\LARGE\bfseries\color{accent} SEPARATION PROCESSES II}\\[2pt]
{\color{accent2} CHE F313 \textperiodcentered{} Problems and Solutions \textperiodcentered{} I-Semester 2026-27}\\[4pt]
{\color{rule}\rule{0.6\linewidth}{0.8pt}}\\[6pt]
{\large\bfseries\color{accent} Every numerical from the uploaded slides, with its solution}\\[3pt]
{\color{accent2} Module 1 Lectures 1--6 \textperiodcentered{} Module 2 Lectures 1--2}
\end{center}

\vspace{2mm}
\begin{flagres}
\textbf{\FLAGTAG{} Read this first.} This document mixes two kinds of statement,
and they are \emph{always} labelled:
\begin{itemize}
  \item \SOURCE{} --- transcribed from the uploaded slide decks in \texttt{Slides/}.
        Questions, given data, formulas and methods. \textbf{These are the truth.}
  \item \DERIVED{} --- algebra, arithmetic and unit conversions performed here.
        \textbf{Generated, so verify before you rely on it.} Nothing in a
        \DERIVED{} box overrides a \SOURCE{} box.
  \item \VERIFY{} --- an independent re-computation used to check an answer.
  \item \TEXT{} --- a textbook value/formula, used \emph{only} as secondary
        confirmation. Cited where used.
  \item \FLAGTAG{} --- an ambiguity, a missing datum, or an assumption. Every
        place where the slides are silent, or where I had to interpret, is
        flagged rather than filled in silently.
\end{itemize}
If a \SOURCE{} box and a \DERIVED{} box ever disagree, the \SOURCE{} box is wrong
and the whole solution must be recomputed. Where a needed datum is absent from
the slides, it is stated as absent --- not invented.
\end{flagres}

\spsub{Question inventory (from the uploaded decks)}
{\small
\begin{tabularx}{\linewidth}{@{}l l X l@{}}
\toprule
\textbf{ID} & \textbf{Source} & \textbf{Topic / what is asked} & \textbf{Solved at}\\
\midrule
Q1 & Mod 1 Lec 3, slide 10 & Clay catalyst: differential \& cumulative analysis, $A_w$, $\overline{D}_s$, particle count & Q1, p.\,\pageref{q:1}\\
Q2 & Mod 1 Lec 3, slides 11--12 (table on 12; also Lec 2 slide 17) & Crushed quartz: $A_w$, $N_w$, $\overline{D}_v$, $\overline{D}_s$, $\overline{D}_w$, $N_i$, number fraction & Q2, p.\,\pageref{q:2}\\
Q3 & Mod 1 Lec 4, slide 19 & Muller mixer: mixing index $I_p$ and standard deviation & Q3, p.\,\pageref{q:3}\\
Q4 & Mod 1 Lec 5, slide 20 & Crusher power by Rittinger's law & Q4, p.\,\pageref{q:4}\\
Q5 & Mod 2 Lec 1, slide 16 (solution graph slide 18) & Quartz screened at 10 mesh: $D/F$, $B/F$, overall effectiveness $E$ & Q5, p.\,\pageref{q:5}\\
\bottomrule
\end{tabularx}}

\notetext{Slide wording is quoted verbatim throughout, \emph{including} the slides' own spelling and small
grammar slips (``It has density of 1.2\,g/cc'', ``analyzed calorimetrically''), so any quotation can be
checked character-for-character against the deck. Two transcription notes. On \texttt{Mod1-Lecture 3} slide 12 the question carries only the
heading \texttt{Numerical Problem}, with the screen-analysis table and the shape factors inside the slide's
image object; on \texttt{Mod1-Lecture 5} slide 20 the powers of ten in \SI{50e-4}{\meter} and friends sit
inside an image object as well. Both were read from those images and are reproduced in the \SOURCE{} boxes
above (the 14-row table was additionally checked row by row, see Q2). Nothing in this document was taken
from a secondary summary of these slides.}

\spsection{REFERENCE NOTE --- WHAT IS TRUSTED, AND IN WHAT ORDER}

\spsub{Precedence when sources disagree}
\begin{enumerate}
  \item \textbf{The uploaded decks} (\texttt{Slides/}, 8 files: \texttt{Mod1-Lecture 1--5.pptx},
        \texttt{Mod1-Lecture 6-compressed.pdf}, \texttt{Mod 2-Lecture 1-compressed.pdf},
        \texttt{Mod 2-Lecture 2.pptx}). This is the course's own material --- the primary source of truth.
  \item \textbf{The cited textbooks} --- McCabe, Smith \& Harriott, \emph{Unit Operations of Chemical
        Engineering}, 7e (Ch.~28, 29), and Swain, Patra \& Roy, \emph{Mechanical Operations}
        (Tata McGraw Hill). Used only to confirm a standard form or a published answer.
  \item \textbf{This document's \DERIVED{} material} --- generated arithmetic. Lowest precedence, always.
\end{enumerate}

\spsub{Rules applied in producing the solutions}
\begin{enumerate}
  \item The \emph{method} used for every question is the method taught on the slides. No outside
        method is substituted.
  \item Every number that appears in a solution came from either a \SOURCE{} box or a computation in
        \texttt{sp2-answers-check.py} (kept alongside this file). No arithmetic was done by hand.
  \item Where the slides fix a required datum (e.g.\ $a=0.8$, $\Phi_s=0.571$, $\rho_p$, $W_i$-style
        inputs), that value is used. Where a datum is \emph{not} on the slides, the solution says so
        and either parameterises it or flags the assumption.
  \item SI/unit conventions follow the slides: $D_p$ in mm with $A_w$ in \si{\square\milli\meter\per\gram}
        for the screen-analysis problems, and \si{\kilo\watt}, \si{\tonne\per\hour} for the crushing problem.
  \item Answers are given with the precision the input data supports, plus an unrounded value where
        useful.
\end{enumerate}

\spsection{Q1 --- SCREEN ANALYSIS OF A CLAY CATALYST}\label{q:1}
\qtitle{Q1}{Module 1 \textperiodcentered{} Lecture 3 \textperiodcentered{} slide 10 \textperiodcentered{} \texttt{Mod1-Lecture 3.pptx}}

\begin{srcres}
\SOURCE{} \textbf{Question as given on the slide} (``Problem 1: Screen Analysis''), quoted verbatim:\\
``Finely divided clay is used as a catalyst in the petroleum industry. It has density of 1.2\,g/cc and a
sphericity of 0.5. The size analysis is as follows:''
\begin{center}\small
\begin{tabular}{@{}l S[table-format=1.4] S[table-format=1.4] S[table-format=1.4] S[table-format=1.4] S[table-format=1.4]@{}}
\toprule
$D_{pi}$ / \si{\milli\meter} & 0.0252 & 0.0178 & 0.0126 & 0.0089 & 0.0038\\
$x_i$ (g/g) & 0.088 & 0.178 & 0.293 & 0.194 & 0.247\\
\bottomrule
\end{tabular}
\end{center}
``Show the differential \& cumulative screen analysis of the mixture. Find the specific surface area and
the Sauter mean diameter of the clay material. Find the number of particles in 1\,g sample.''
\end{srcres}

\begin{srcres}
\SOURCE{} \textbf{Method from the slides} (Mod 1 Lec 3, derivations 1--13):
\[
A_w=\frac{6}{\rho_p\Phi_s}\sum_{i=1}^{n}\frac{x_i}{D_{pi}},\qquad
\overline{D}_s=\frac{1}{\sum_i x_i/D_{pi}}=\frac{6}{\Phi_s\rho_p A_w},\qquad
N_w=\frac{1}{a\rho_p}\sum_{i=1}^{n}\frac{x_i}{D_{pi}^3}
\]
Cumulative analysis: add the increments consecutively starting from the smallest particles, and plot
the sums against the maximum particle diameter of the increment (Lec 2 slide 16).
\end{srcres}

\spsub{Step 1 --- differential and cumulative analysis}
\begin{devres}
\DERIVED{} The differential analysis is the tabulated $x_i$ versus $D_{pi}$ already given. The cumulative
``fraction smaller than $D_{pi}$'' is the running sum from the smallest size; ``fraction larger'' is its
complement. Check of the given data: $\sum x_i = 0.088+0.178+0.293+0.194+0.247 = 1.000$.
\end{devres}
{\small
\begin{tabular}{@{}S[table-format=1.4] S[table-format=1.3] S[table-format=1.4] S[table-format=1.4]@{}}
\toprule
{$D_{pi}$ / \si{\milli\meter}} & {$x_i$ (differential)} & {cumulative smaller} & {cumulative larger}\\
\midrule
0.0038 & 0.247 & 0.2470 & 0.7530\\
0.0089 & 0.194 & 0.4410 & 0.5590\\
0.0126 & 0.293 & 0.7340 & 0.2660\\
0.0178 & 0.178 & 0.9120 & 0.0880\\
0.0252 & 0.088 & 1.0000 & 0.0000\\
\bottomrule
\end{tabular}}
\notetext{\DERIVED{} Rows are listed from the smallest increment upward, which is the order used when building
a cumulative distribution.}

\spsub{Step 2 --- specific surface area and Sauter mean diameter}
\begin{devres}
\DERIVED{} Unit note: $D_{pi}$ is in \si{\milli\meter}, so $\rho_p$ must be in \si{\gram\per\cubic\milli\meter}:
\[
\rho_p=\SI{1.2}{\gram\per\cubic\centi\meter}=\SI{1.2e-3}{\gram\per\cubic\milli\meter}
\]
\[
\sum_i\frac{x_i}{D_{pi}}=123.544\ \text{mm}^{-1}
\qquad(\text{e.g. } 0.088/0.0252=3.4921,\ 0.247/0.0038=65.000)
\]
\[
A_w=\frac{6}{\rho_p\Phi_s}\sum_i\frac{x_i}{D_{pi}}
=\frac{6}{(\num{1.2e-3})(0.5)}(123.544)
=\SI{1.2354e6}{\square\milli\meter\per\gram}
\]
\[
\overline{D}_s=\frac{1}{\sum_i x_i/D_{pi}}=\frac{1}{123.544}
=\SI{0.008094}{\milli\meter}=\SI{8.09}{\micro\meter}
\]
\end{devres}
\begin{verres}
\VERIFY{} Consistency of the two slide formulas: $\dfrac{6}{\Phi_s\rho_p A_w}
=\dfrac{6}{(0.5)(\num{1.2e-3})(\num{1.2354e6})}=\SI{0.008094}{\milli\meter}=\overline{D}_s$ \checkmark\\[2pt]
Sanity: $A_w=\SI{1.2354e6}{\square\milli\meter\per\gram}=\SI{1.24}{\square\meter\per\gram}$, which is
the right order for a powder of $\sim\SI{8}{\micro\meter}$ particles.
\end{verres}

\spsub{Step 3 --- number of particles in a \SI{1}{\gram} sample}
\begin{devres}
\DERIVED{} For the same data,
\[
\sum_i\frac{x_i}{D_{pi}^3}=4.9601\times10^{6}\ \text{mm}^{-3},
\qquad
N_w=\frac{1}{a\rho_p}\sum_i\frac{x_i}{D_{pi}^3}
=\frac{4.9601\times10^{6}}{a\,(\num{1.2e-3})}
=\frac{4.1334\times10^{9}}{a}\ \text{g}^{-1}
\]
The slide does not state the volume shape factor $a$, so the result is left in terms of it. For the
sphere value used elsewhere in the course ($a=\pi/6=0.5236$, Lec 3 slide 6 and the textbook),
\[
N_w=\frac{4.1334\times10^{9}}{0.5236}=\SI{7.89e9}{particles\per\gram}
\]
\end{devres}
\begin{flagres}
\FLAGTAG{} $a$ is \emph{not} given on this slide (unlike Q2, where $a=0.8$ is stated). The value
$\SI{7.9e9}{\per\gram}$ assumes $a=\pi/6$ (spheres). Since $N_w\propto 1/a$, quoting $a$ from the
question sheet changes the answer: $a=0.5\to\SI{8.3e9}{\per\gram}$, $a=0.8\to\SI{5.2e9}{\per\gram}$.
Confirm the intended $a$ before submitting.
\end{flagres}

\begin{ansbox}
\textbf{Answers (Q1).}\\
(a) Differential and cumulative tables as above ($\sum x_i=1.000$).\\
(b) $A_w=\SI{1.2354e6}{\square\milli\meter\per\gram}\;(\approx\SI{1.24}{\square\meter\per\gram})$;\qquad
    $\overline{D}_s=\SI{0.00809}{\milli\meter}=\SI{8.09}{\micro\meter}$.\\
(c) $N_w=(4.133\times10^{9})/a$ particles per gram; with $a=\pi/6$:
    $\approx\SI{7.9e9}{particles\per\gram}$. \emph{($a$ not given on the slide --- see the flag.)}
\end{ansbox}

\newpage
\spsection{Q2 --- SCREEN ANALYSIS OF CRUSHED QUARTZ}\label{q:2}
\qtitle{Q2}{Module 1 \textperiodcentered{} Lecture 3 \textperiodcentered{} slides 11--12 (table on slide 12) \textperiodcentered{} also Lec 2 slide 17}

\begin{srcres}
\SOURCE{} \textbf{Question as given on the slide} (``Problem 2: Screen Analysis''), quoted verbatim:\\
``The screen analysis shown in Table applies to a sample of crushed quartz. The density of the particles
is 2650\,kg/m3 (0.00265\,g/mm3), and the shape factors are a = 0.8 and $\Phi_s = 0.571$. For the material
between 4-mesh and 200-mesh in particle size, calculate (a) Aw in square millimeters per gram and Nw in
particles per gram, (b) Dv (c) Ds (d) Dw and (e) Ni for the 150/200-mesh increment.
(f) What fraction of the total number of particles is in the 150/200-mesh increment?''
\end{srcres}

\begin{srcres}
\SOURCE{} \textbf{Data table as printed on slide 12} (the slide's only text is the heading
\texttt{Numerical Problem}; the same table also appears on Lecture 2 slide 17). Columns as on the slide:
mesh, screen opening, mass fraction retained $x_i$, average particle diameter in increment, cumulative
fraction smaller than $D_{pi}$.
\begin{center}\footnotesize
\begin{tabular}{@{}l S[table-format=1.3] S[table-format=1.4] S[table-format=1.3] S[table-format=1.4]@{}}
\toprule
{Mesh} & {Screen opening $D_{pi}$/\si{\milli\meter}} & {retained $x_i$} & {$D_{pi}$ in increment/\si{\milli\meter}} & {cumulative smaller}\\
\midrule
4   & 4.699 & 0.0000 & {---}   & 1.0000\\
6   & 3.327 & 0.0251 & 4.013 & 0.9749\\
8   & 2.362 & 0.1250 & 2.845 & 0.8499\\
10  & 1.651 & 0.3207 & 2.007 & 0.5292\\
14  & 1.168 & 0.2570 & 1.409 & 0.2722\\
20  & 0.833 & 0.1590 & 1.001 & 0.1132\\
28  & 0.589 & 0.0538 & 0.711 & 0.0594\\
35  & 0.417 & 0.0210 & 0.503 & 0.0384\\
48  & 0.295 & 0.0102 & 0.356 & 0.0282\\
65  & 0.208 & 0.0077 & 0.252 & 0.0205\\
100 & 0.147 & 0.0058 & 0.178 & 0.0147\\
150 & 0.104 & 0.0041 & 0.126 & 0.0106\\
200 & 0.074 & 0.0031 & 0.089 & 0.0075\\
Pan & {---}   & 0.0075 & 0.037 & 0.0000\\
\bottomrule
\end{tabular}
\end{center}
\notetext{\DERIVED{} Internal check of the table: each increment diameter is the arithmetic mean of the two
defining screen openings, e.g.\ $(4.699+3.327)/2=4.013$, $(3.327+2.362)/2=2.845$, \dots,
$(0.074+0)/2=0.037$ for the pan. All 13 values reproduce the table to within \SI{0.0005}{\milli\meter}
(the table prints three decimals, so e.g.\ the 14-mesh mean $1.4095$ is printed as $1.409$); the data is
therefore self-consistent and no reading error is involved.
The 12 increments from 4-mesh to 200-mesh carry $\sum x_i=0.9925$; the remaining $0.0075$ is the pan.
Because the question says ``between 4-mesh and 200-mesh'', the pan is \emph{excluded} from the answer
(the textbook solution does the same) --- see the flag below for the numbers if it is included.}
\end{srcres}

\begin{srcres}
\SOURCE{} \textbf{Method from the slides}:
\[
A_w=\frac{6}{\rho_p\Phi_s}\sum\frac{x_i}{D_{pi}},\qquad
N_w=\frac{1}{a\rho_p}\sum\frac{x_i}{D_{pi}^3},\qquad
\overline{D}_v=\Big[\frac{1}{\sum x_i/D_{pi}^3}\Big]^{1/3},
\]
\[
\overline{D}_s=\frac{1}{\sum x_i/D_{pi}},\qquad
\overline{D}_w=\sum x_i D_{pi},\qquad
N_i=\frac{m_i}{a\rho_p D_{pi}^3}=\frac{x_i}{a\rho_p\overline{D}_{pi}^3}\ \text{(per unit mass)}
\]
\end{srcres}

\spsub{Step 1 --- the two sums (12 increments, 4-mesh to 200-mesh)}
\begin{devres}
\DERIVED{}
\[
\sum_{i=1}^{12}\frac{x_i}{D_{pi}}=0.82780\ \text{mm}^{-1}
\qquad\qquad
\sum_{i=1}^{12}\frac{x_i}{D_{pi}^3}=8.7932\ \text{mm}^{-3}
\]
work for the first increment, $D_{pi}=4.013$ and $x_i=0.0251$:
$\;0.0251/4.013=0.00626$ and $0.0251/4.013^3=0.000389$.
\end{devres}

\spsub{Step 2 --- (a) specific surface area and number of particles}
\begin{devres}
\DERIVED{}
\[
A_w=\frac{6}{\rho_p\Phi_s}\sum\frac{x_i}{D_{pi}}
=\frac{6}{(\num{2.65e-3})(0.571)}(0.82780)
=\SI{3282}{\square\milli\meter\per\gram}=\SI{32.8}{\square\centi\meter\per\gram}
\]
\[
N_w=\frac{1}{a\rho_p}\sum\frac{x_i}{D_{pi}^3}
=\frac{8.7932}{(0.8)(\num{2.65e-3})}
=\SI{4148}{particles\per\gram}
\]
\end{devres}

\spsub{Step 3 --- (b) volume mean, (c) Sauter, (d) mass mean diameters}
\begin{devres}
\DERIVED{}
\[
\overline{D}_v=\Big[\frac{1}{\sum x_i/D_{pi}^3}\Big]^{1/3}
=\Big[\frac{1}{8.7932}\Big]^{1/3}=\SI{0.4845}{\milli\meter}
\]
\[
\overline{D}_s=\frac{1}{\sum x_i/D_{pi}}=\frac{1}{0.82780}=\SI{1.2080}{\milli\meter}
\]
\[
\overline{D}_w=\sum x_i D_{pi}
=0.0251(4.013)+0.1250(2.845)+\dots+0.0031(0.089)=\SI{1.6775}{\milli\meter}
\]
\end{devres}

\spsub{Step 4 --- (e) particles in the 150/200-mesh increment}
\begin{devres}
\DERIVED{} For that increment $x_i=0.0031$ and $D_{pi}=\SI{0.089}{\milli\meter}$ (per unit mass of the
whole sample):
\[
N_i=\frac{x_i}{a\rho_p D_{pi}^3}
=\frac{0.0031}{(0.8)(\num{2.65e-3})(0.089)^3}
=\SI{2074}{particles\per\gram}
\]
\end{devres}

\spsub{Step 5 --- (f) fraction of the total number of particles}
\begin{devres}
\DERIVED{} The fraction of all particles that lie in that increment is its share of
$\sum x_i/D_{pi}^3$:
\[
\frac{N_i}{N_T}=\frac{x_i/D_{pi}^3}{\sum_j x_j/D_{pj}^3}
=\frac{0.0031/0.089^3}{8.7932}=0.5001\ \approx\ 50\%
\]
Note the contrast with its \emph{mass} fraction, $x_i=0.0031=0.31\%$: half of all the particles sit in
this one fine increment even though it is a negligible part of the mass.
\end{devres}

\begin{verres}
\VERIFY{} Independent checks: (i) the algebra route
$N_i/N_T=\big[(1/(a\rho_p))\,(x_i/D_{pi}^3)\big]/N_w=2074/4148=0.500$ agrees;
(ii) the published worked solution of this textbook problem, as reproduced in course material,
gives $\sum x_i/D_{pi}=0.8278$, $\sum x_i/D_{pi}^3=8.7932$,
$A_w=\SI{3282}{\square\milli\meter\per\gram}$, $N_w=\SI{4148}{particles\per\gram}$,
$N_i=2074$ and a fraction of $50\%$ \TEXT{} --- i.e.\ this solution reproduces the published
answer to the digits quoted, which is strong evidence that the table was transcribed correctly.
\end{verres}

\begin{flagres}
\FLAGTAG{} Two points to confirm with the instructor:
\begin{enumerate}
  \item \textbf{Scope.} ``Between 4-mesh and 200-mesh'' excludes the pan, so $\sum x_i=0.9925$ and the
        denominators are not normalised to 1. Including the pan gives different numbers
        ($A_w=\SI{4086}{\square\milli\meter\per\gram}$, $N_w=\SI{73990}{particles\per\gram}$,
        $\overline{D}_v=\SI{0.185}{\milli\meter}$, $\overline{D}_s=\SI{0.970}{\milli\meter}$,
        $\overline{D}_w=\SI{1.678}{\milli\meter}$, $N_i/N_T=2.8\%$). The exclusion is what the
        textbook solution does, so it is used above.
  \item \textbf{Which diameter to use.} The solution uses the increment-mean column $D_{pi}$ from the
        table. The 150/200 value there is $0.089$\,mm (not $0.104$ or $0.074$, the screen openings).
\end{enumerate}
\end{flagres}

\begin{ansbox}
\textbf{Answers (Q2).}\\
(a) $A_w=\SI{3282}{\square\milli\meter\per\gram}$, \quad $N_w=\SI{4148}{particles\per\gram}$.\\
(b) $\overline{D}_v=\SI{0.484}{\milli\meter}$.\\
(c) $\overline{D}_s=\SI{1.208}{\milli\meter}$.\\
(d) $\overline{D}_w=\SI{1.677}{\milli\meter}$.\\
(e) $N_i=\SI{2074}{particles\per\gram}$ \;(150/200-mesh increment).\\
(f) $N_i/N_T=0.500$ --- about \textbf{50\%} of all the particles, against only $0.31\%$ of the mass.
\end{ansbox}

\newpage
\spsection{Q3 --- MIXING INDEX IN A MULLER MIXER}\label{q:3}
\qtitle{Q3}{Module 1 \textperiodcentered{} Lecture 4 \textperiodcentered{} slide 19 \textperiodcentered{} \texttt{Mod1-Lecture 4.pptx}}

\begin{srcres}
\SOURCE{} \textbf{Question as given on the slide} (``Problem 1''), quoted verbatim:\\
``A silty soil containing 14 percent moisture was mixed in a large muller mixer with 10\,wt\% of a tracer
consisting of dextrose and picric acid. After 3\,min of mixing, 12 random samples were taken from the mix
and analyzed calorimetrically for tracer material. The measured concentrations in the sample were in
weight percent tracer, 10.24, 9.30, 7.94, 10.24, 11.08, 10.03, 11.91, 9.72, 9.20, 10.76, 10.97, 10.55.
Calculate the mixing index IP and the standard deviation.''
\end{srcres}

\begin{srcres}
\SOURCE{} \textbf{Method from the slides} (Mod 1 Lec 4, cohesive solids / pastes --- and Lec 4 slide 16 for $s$):
\[
s=\sqrt{\frac{\sum_{i=1}^{N}(x_i-\bar{x})^2}{N-1}}
=\sqrt{\frac{\sum_{i=1}^{N}x_i^2-\bar{x}\sum_{i=1}^{N}x_i}{N-1}},
\qquad
\sigma_0=\sqrt{\mu(1-\mu)},
\qquad
I_s=\frac{\sigma_0}{s}
\]
``For cohesive solids, instead of number, mass fraction is used''; $x_i$ = mass fraction of A in each
sample, $\bar x$ = average of the mass fractions, $N$ = number of spot samples; $\sigma_0$ is the
theoretical standard deviation at zero mixing. The soil is wet and cohesive, and the samples were
analysed for tracer \emph{mass} percent, so the paste equations apply.
\end{srcres}

\spsub{Step 1 --- standard deviation of the samples}
\begin{devres}
\DERIVED{}
\[
\bar{x}=\frac{1}{12}\sum_{i=1}^{12}x_i=\frac{121.94}{12}=10.1617\ \text{wt\%}
\qquad
\sum_{i=1}^{12}(x_i-\bar{x})^2=11.890\ (\text{wt\%})^2
\]
\[
s=\sqrt{\frac{11.890}{12-1}}=\SI{1.0397}{wt\%}=0.010397\ \text{(mass fraction)}
\]
(The $N$ instead of $N-1$ denominator would give $s=\SI{0.9954}{wt\%}=\SI{0.009954}{}$ as a fraction.)
\end{devres}

\spsub{Step 2 --- theoretical standard deviation at zero mixing}
\begin{devres}
\DERIVED{} The mixture was charged with \SI{10}{wt\%} tracer, so $\mu=0.10$:
\[
\sigma_0=\sqrt{\mu(1-\mu)}=\sqrt{0.10\times0.90}=\sqrt{0.09}=0.30
\]
\end{devres}

\spsub{Step 3 --- mixing index}
\begin{devres}
\DERIVED{}
\[
I_p=\frac{\sigma_0}{s}=\frac{0.30}{0.010397}=28.9
\]
Since $s$ has fallen well below $\sigma_0$, the value of $I_p$ strongly above 1 indicates that the
mixer has mixed the tracer far more uniformly than a completely unmixed, segregated charge would be
(compare $I_s=\sigma_e/s$ for granular solids, which rises from 0 towards 1 as mixing proceeds and is
capped near 1 for a random mixture).
\end{devres}

\begin{flagres}
\FLAGTAG{} Four judgement calls, all of which change the last digit or more:
\begin{enumerate}
  \item \textbf{$\mu$: charged or measured?} $\mu=0.10$ is used, because that is the tracer actually
        charged (``mixed with 10 wt\% of a tracer''). Using the measured mean instead,
        $\mu=\bar x=0.1016$, gives $\sigma_0=0.3022$ and $I_p=29.1$.
  \item \textbf{$N-1$ or $N$?} The slide's formula is $s=\sqrt{\sum(x_i-\bar{x})^2/(N-1)}$, used above.
        Some textbooks use $N$: that gives $s=\SI{0.995}{\percent}$ and $I_p=30.1$.
  \item \textbf{Number vs mass fraction.} The question says ``silty soil'' (cohesive, 14\% moisture)
        and reports wt\% tracer, so the paste form $\sigma_0=\sqrt{\mu(1-\mu)}$ with mass fractions is
        used. The granular form $\sigma_e=\sqrt{\mu_p(1-\mu_p)/n}$ needs $n$, the number of particles
        per sample, which the slide does not give.
  \item \textbf{Symbol.} The slide writes the paste mixing index as $I_s$; the question calls it
        $I_p$. They are the same quantity.
\end{enumerate}
\end{flagres}

\begin{ansbox}
\textbf{Answers (Q3).}\\
$\bar x = 10.16$\,wt\%, \quad $s=\SI{1.04}{wt\%}$ (i.e.\ $0.0104$ as a mass fraction; $s=0.995$\,wt\% if
the $N$ denominator is used), \quad $\sigma_0=0.30$, \quad
$I_p=0.30/0.0104=28.9$ \;($\approx 29$).\\
With the measured mean as $\mu$ and the $N$-denominator, the same working gives $I_p\approx30$.
\end{ansbox}

\newpage
\spsection{Q4 --- CRUSHER POWER BY RITTINGER'S LAW}\label{q:4}
\qtitle{Q4}{Module 1 \textperiodcentered{} Lecture 5 \textperiodcentered{} slide 20 \textperiodcentered{} \texttt{Mod1-Lecture 5.pptx}}

\begin{srcres}
\SOURCE{} \textbf{Question as given on the slide} (``Problem 1: Size reduction''), quoted verbatim:\\
``Particles of the average feed size of $50\times10^{-4}$\,m are crushed to an average product size of
$10\times10^{-4}$\,m at the rate of 20 tonnes per hour. At this rate, the crusher consumes
\SI{40}{\kilo\watt} of power which \SI{5}{\kilo\watt} are required for running the mill empty.
Calculate the power consumption if 12 tonnes/h of this product is crushed to
$5\times10^{-4}$\,m size in the same mill? Assume Rittinger's Law is applicable.''
\end{srcres}

\begin{srcres}
\SOURCE{} \textbf{Method fixed by the slide} (Mod 1 Lec 5, Rittinger's law):
\[
\frac{P}{\dot m}=K_R\Big[\frac{1}{\overline{D}_{sb}}-\frac{1}{\overline{D}_{sa}}\Big]
\]
with $\overline{D}_{sa}$ = feed (a) size and $\overline{D}_{sb}$ = product (b) size. \emph{The question
directs Rittinger's law, so it is used --- not Kick's or Bond's.}
\end{srcres}

\spsub{Step 1 --- useful (net) power of the first duty}
\begin{devres}
\DERIVED{} The \SI{5}{\kilo\watt} is spent running the mill empty, so the power actually doing crushing is
\[
P_{\text{net},1}=40-5=\SI{35}{\kilo\watt}\ \text{at}\ \dot m_1=\SI{20}{\tonne\per\hour}
\]
\[
\frac{P_{\text{net},1}}{\dot m_1}=\frac{35}{20}=\SI{1.75}{\kilo\watt\hour\per\tonne}
\]
\end{devres}

\spsub{Step 2 --- the Rittinger constant of this mill and material}
\begin{devres}
\DERIVED{} With $D_{sa}=50\times10^{-4}$\,m and $D_{sb}=10\times10^{-4}$\,m,
\[
\frac{1}{D_{sb}}-\frac{1}{D_{sa}}=\frac{1}{10\times10^{-4}}-\frac{1}{50\times10^{-4}}
=1000-200=\SI{800}{\per\meter}
\]
\[
K_R=\frac{P/\dot m}{\big[1/D_{sb}-1/D_{sa}\big]}=\frac{1.75}{800}
=\SI{2.1875e-3}{\kilo\watt\hour\meter\per\tonne}
\]
\end{devres}

\spsub{Step 3 --- the second duty}
\begin{devres}
\DERIVED{} Now the same mill crushes \emph{this product} ($D_{sa}=10\times10^{-4}$\,m) down to
$5\times10^{-4}$\,m at $\dot m_2=\SI{12}{\tonne\per\hour}$:
\[
\frac{1}{D_{sb}}-\frac{1}{D_{sa}}=\frac{1}{5\times10^{-4}}-\frac{1}{10\times10^{-4}}
=2000-1000=\SI{1000}{\per\meter}
\]
\[
P_{\text{net},2}=K_R\,\dot m_2\Big[\frac{1}{D_{sb}}-\frac{1}{D_{sa}}\Big]
=(\num{2.1875e-3})(12)(1000)=\SI{26.25}{\kilo\watt}
\]
Restoring the fixed empty-mill load:
\[
P_2=P_{\text{net},2}+5=\SI{31.25}{\kilo\watt}
\]
\end{devres}

\spsub{Step 4 --- the same calculation done entirely with power ratios}
\begin{devres}
\DERIVED{} A second route, keeping the mill load in the ratio:
\[
\frac{P_2-5}{P_1-5}
=\frac{\dot m_2}{\dot m_1}\cdot
\frac{\big[1/D_{sb,2}-1/D_{sa,2}\big]}{\big[1/D_{sb,1}-1/D_{sa,1}\big]}
=\frac{12}{20}\cdot\frac{1000}{800}=0.75
\]
\[
P_2=5+0.75(40-5)=5+26.25=\SI{31.25}{\kilo\watt}
\]
\end{devres}

\begin{verres}
\VERIFY{} The two routes agree exactly ($26.25+5=31.25$), and both are dimensionally consistent:
\si{\kilo\watt\hour\per\tonne} $\times$ \si{\tonne\per\hour} $\times$ \si{\per\meter} must be checked
with $D$ in metres, giving \si{\kilo\watt}. Note also that the second duty crushes a \emph{finer} feed
($10\to5$ instead of $50\to10$), and the size-reduction step
$(1/5-1/10)=1000$ exceeds the first duty's $(1/10-1/50)=800$ per metre, so the smaller throughput does
not make the power small.
\end{verres}

\begin{flagres}
\FLAGTAG{} Two readings of ``power consumption'', both shown:
\begin{itemize}
  \item \textbf{Total power drawn by the crusher: \SI{31.25}{\kilo\watt}} --- the answer above, because
        the question asks for the \emph{power consumption} of the machine and the empty-mill load is
        still drawn.
  \item \textbf{Net crushing power: \SI{26.25}{\kilo\watt}} --- if only the work of size reduction is
        wanted. Quoting \SI{26.25}{\kilo\watt} as the machine's consumption would double-count the
        treatment of the \SI{5}{\kilo\watt} already handled in the first duty.
\end{itemize}
\end{flagres}

\begin{ansbox}
\textbf{Answers (Q4).} $K_R=\SI{2.1875e-3}{\kilo\watt\hour\meter\per\tonne}$;
$P_{\text{net},2}=\SI{26.25}{\kilo\watt}$; \quad
$\boxed{P_2=\SI{31.25}{\kilo\watt}}$ including the empty-mill power.
\end{ansbox}

\newpage
\spsection{Q5 --- SCREEN EFFECTIVENESS OF A QUARTZ MIXTURE}\label{q:5}
\qtitle{Q5}{Module 2 \textperiodcentered{} Lecture 1 \textperiodcentered{} slide 16, solution graph slide 18 \textperiodcentered{} \texttt{Mod 2-Lecture 1-compressed.pdf}}

\begin{srcres}
\SOURCE{} \textbf{Question as given on the slide} (``Problem on Screen Effectiveness''), quoted verbatim:\\
``A quartz mixture having the screen analysis shown in Table is screened through a standard 10-mesh
screen. The cumulative screen analysis of overflow and underflow are given in Table. Calculate --- the
mass ratios of the overflow and underflow to feed and --- the overall effectiveness of the screen.''
\end{srcres}

\begin{srcres}
\SOURCE{} \textbf{Data table as printed on slide 16} --- the slide's caption is ``Screen analyses ---
Cumulative fraction LARGER than $D_p$'', with columns Feed, Overflow, Underflow:
\begin{center}\small
\begin{tabular}{@{}l S[table-format=1.3] S[table-format=1.3] S[table-format=1.2] S[table-format=1.3]@{}}
\toprule
{Mesh} & {$D_p$/\si{\milli\meter}} & {Feed $x_F$} & {Overflow $x_D$} & {Underflow $x_B$}\\
\midrule
4  & 4.699 & 0     & 0     & \\
6  & 3.327 & 0.025 & 0.071 & \\
8  & 2.362 & 0.15  & 0.43  & 0\\
10 & 1.651 & 0.47  & 0.85  & 0.195\\
14 & 1.168 & 0.73  & 0.97  & 0.58\\
20 & 0.833 & 0.885 & 0.99  & 0.83\\
28 & 0.589 & 0.94  & 1.00  & 0.91\\
35 & 0.417 & 0.96  &       & 0.94\\
65 & 0.208 & 0.98  &       & 0.975\\
Pan&       & 1.00  &       & 1.00\\
\bottomrule
\end{tabular}
\end{center}
\end{srcres}

\notetext{\DERIVED{} Reading of the data: the cut diameter is the 10-mesh opening,
$D_{pc}=\SI{1.651}{\milli\meter}$ (the slide says the test screen is 10-mesh), so the table row for
10 mesh supplies the three values needed directly --- $x_F=0.47$, $x_D=0.85$, $x_B=0.195$. No
interpolation is needed, because the slide's table is already tabulated at that size. Blank cells in the
overflow column (35 and 65 mesh) mean that no overflow stream was analysed above 28 mesh; they are not
needed for the calculation, which uses only the 10-mesh row.}

\begin{srcres}
\SOURCE{} \textbf{Method from the slides} (Mod 2 Lec 1, hand-worked slides 12--15):\\
Overall and oversize balances give the split,
\[
\frac{D}{F}=\frac{x_F-x_B}{x_D-x_B},\qquad
\frac{B}{F}=\frac{x_F-x_D}{x_B-x_D}
\]
and the effectiveness of the screen in each direction,
\[
E_A=\frac{D x_D}{F x_F},\qquad
E_B=\frac{B(1-x_B)}{F(1-x_F)},\qquad
E=E_A E_B
\]
with the rearranged form
$E=\dfrac{x_D}{x_F}\dfrac{x_F-x_B}{x_D-x_B}\Big[1-\dfrac{1-x_D}{1-x_F}\dfrac{x_F-x_B}{x_D-x_B}\Big]$.
The slide's own worked solution is the graphical one on slide 18: the cumulative-fraction curves for
feed, overflow and underflow are plotted together against $D_p$, and read off at the cut diameter.
\end{srcres}

\spsub{Step 1 --- mass ratios of overflow and underflow to feed}
\begin{devres}
\DERIVED{}
\[
\frac{D}{F}=\frac{x_F-x_B}{x_D-x_B}=\frac{0.47-0.195}{0.85-0.195}=\frac{0.275}{0.655}=0.4198
\]
\[
\frac{B}{F}=\frac{x_F-x_D}{x_B-x_D}=\frac{0.47-0.85}{0.195-0.85}=\frac{-0.38}{-0.655}=0.5802
\]
Mass-balance closure: $D/F+B/F=0.4198+0.5802=1.000$ \checkmark
\end{devres}

\spsub{Step 2 --- effectiveness of each fraction}
\begin{devres}
\DERIVED{}
\[
E_A=\frac{D}{F}\cdot\frac{x_D}{x_F}=0.4198\times\frac{0.85}{0.47}=0.4198\times1.8085=0.7593
\]
\[
E_B=\frac{B}{F}\cdot\frac{1-x_B}{1-x_F}=0.5802\times\frac{0.805}{0.53}=0.5802\times1.5189=0.8812
\]
\end{devres}

\spsub{Step 3 --- overall effectiveness}
\begin{devres}
\DERIVED{}
\[
E=E_A E_B=0.7593\times0.8812=0.6691\ \approx\ 66.9\%
\]
The same result from the slide's boxed form:
\[
E=\frac{0.85}{0.47}\cdot\frac{0.275}{0.655}\Big[1-\frac{0.15}{0.53}\cdot\frac{0.275}{0.655}\Big]
=1.8085\times0.4198\times(1-0.1188)=0.6691
\]
\end{devres}

\begin{verres}
\VERIFY{} Both the $E_A E_B$ product and the rearranged boxed form give $0.6691$; the component
balances also close:
\[
\frac{D}{F}x_D+\frac{B}{F}x_B=0.4198(0.85)+0.5802(0.195)=0.470=x_F\ \checkmark
\]
\[
\frac{D}{F}(1-x_D)+\frac{B}{F}(1-x_B)=0.4198(0.15)+0.5802(0.805)=0.530=1-x_F\ \checkmark
\]
A 10-mesh screen recovering $76\%$ of the oversize and $88\%$ of the undersize is a healthy result for a
commercial screen --- consistent with the slide's remark that a commercial screen gives poor separation
compared with a testing screen.
\end{verres}

\begin{flagres}
\FLAGTAG{} Three notes:
\begin{enumerate}
  \item The solution assumes the table is to be read exactly \emph{at} the cut size. Strictly, $x_F$,
        $x_D$ and $x_B$ should be evaluated on the cumulative curves at $D_p=D_{pc}=\SI{1.651}{\milli\meter}$;
        the table row at 10 mesh provides those values directly, so no interpolation arises here.
  \item The underflow column is blank above 8 mesh (nothing coarse reports to the underflow), which is
        consistent with $x_B=0$ at 8 mesh rather than missing data.
  \item \textbf{The boxed rearrangement on the slide.} Slides 12--15 carry a hand-written derivation
        in which the constant term of the rearranged expression is scribbled over. This does not affect
        the answer, because the slides \emph{define} the overall effectiveness as the product of the
        two direction efficiencies, $E=E_AE_B$; the boxed form evaluates to exactly the same number
        ($0.6691$), as verified above, and it also returns the correct limit $E=1$ for a perfect
        separation ($x_D=1$, $x_B=0$). Both checks are in \texttt{sp2-answers-check.py}.
\end{enumerate}
\end{flagres}

\begin{ansbox}
\textbf{Answers (Q5).}\\
$\dfrac{D}{F}=0.42$ (overflow/feed), \quad $\dfrac{B}{F}=0.58$ (underflow/feed).\\
$E_A=0.759$, \quad $E_B=0.881$, \quad
$\boxed{E=E_A E_B = 0.669 \approx 67\%}$ overall effectiveness.
\end{ansbox}

\newpage
\spsection{APPENDIX --- HOW THESE ANSWERS WERE VERIFIED}

\spsub{Reproducibility}
Every \DERIVED{} number in this document was produced by \texttt{sp2-answers-check.py}, kept beside this
file. Running it prints each quantity in the same order as the solutions:
\begin{quote}\small
\texttt{python3 sp2-answers-check.py}
\end{quote}
No arithmetic in the solutions was done by hand, so any value can be re-checked by re-running that script
or by editing the input data at the top of each block.

\spsub{What was checked, and against what}
{\small
\begin{tabularx}{\linewidth}{@{}l X l@{}}
\toprule
\textbf{Item} & \textbf{Check performed} & \textbf{Result}\\
\midrule
Q1 & $\sum x_i=1$; the two independent slide forms for $\overline{D}_s$ agree & pass\\
Q2 & every increment diameter equals the mean of its two defining screen openings & pass (13/13, to \SI{0.0005}{\milli\meter})\\
Q2 & the two routes to $N_i/N_T$ agree; comparison with the published worked solution of this problem & pass (all values)\\
Q3 & slide formula for $s$ read from the slide (image, not OCR); both denominators reported & pass\\
Q4 & ratio route reproduces the direct route; dimensionality of $K_R$ & pass\\
Q5 & $E=E_AE_B$ and the boxed rearranged form agree ($\Delta<10^{-6}$); both balances close; perfect-separation limit gives $E=1$ & pass\\
\bottomrule
\end{tabularx}}

\spsub{Known limits of this document}
\begin{itemize}
  \item It contains the \emph{five} numericals present in the eight uploaded decks. Slides that merely
        show a demonstration video, a diagram, or a worked graph were checked and are not treated as
        questions (see the exclusion list below); no question was invented.
  \item Where the slides give no method --- they never do for the five above --- the method was taken
        from the lecture that set the question, never from outside.
  \item Answers were re-derived from the data, so if an instructor's answer differs, compare the
        flagged assumptions first; they are the only places where the sources leave room.
\end{itemize}

\spsub{Slides examined and excluded (not questions)}
{\footnotesize
\begin{tabularx}{\linewidth}{@{}l X@{}}
\toprule
\textbf{Slide} & \textbf{Why excluded}\\
\midrule
Mod 1 Lec 5, slides 7--18 & hand-written derivations and diagrams, no question posed\\
Mod 1 Lec 3, slides 3--8 & hand-written derivation of average diameters, no question posed\\
Mod 2 Lec 1, slides 11--15 & hand-written derivation of screen effectiveness (source of Q5's method)\\
Mod 1 Lec 2, slide 13; Lec 1 slide 13; Lec 5 slide 17 & data tables/graphs shown as illustration\\
Mod 2 Lec 2, slides 5--12, 20 & diagrams and pressure-gradient profile\\
Various slides with \texttt{youtu.be}/\texttt{youtube.com} links & demonstration videos\\
\bottomrule
\end{tabularx}}

\vspace{3mm}
{\color{rule}\hrule height 0.6pt}
\vspace{2mm}
\notetext{\textbf{Provenance summary.} Questions, data, formulas and methods: \SOURCE{} the uploaded decks.
Published-answer confirmation for Q2: \TEXT{} the crushed-quartz screen-analysis example of McCabe, Smith
\& Harriott, \emph{Unit Operations of Chemical Engineering}, Ch.~28 (consulted as a cross-check of the
numbers reported above; the method itself is taken from the slides). Everything else --- algebra, arithmetic, unit
conversions, the checks above: \DERIVED{} / \VERIFY{} generated. Ambiguities: \FLAGTAG{} flagged in place.
This file is the editable source; regenerate the PDF after any change.}

\end{document}
